\begin{table}[!htbp]\centering
\small
\caption{Absorption of the January-2024 fraud escalation: difference-in-differences estimates}
\label{tab:did}
\resizebox{\textwidth}{!}{%
\begin{tabular}{lrrrrrrr}
\toprule
Specification & Fintech x Post & (CR1 s.e.) & p (CR1) & p (WCR) & p (perm.) & Entities & Fintechs \\
\midrule
(1) All fintech vs banks & -1.059*** & (0.207) & 0.000 & 0.009 & 0.005 & 30 & 5 \\
(2) + lagged ln clients & -0.983*** & (0.247) & 0.000 & 0.015 & 0.009 & 30 & 5 \\
(3) vs microfinance banks & -0.271 & (0.232) & 0.277 & 0.291 & 0.309 & 9 & 5 \\
(4) vs large banks & -0.983*** & (0.241) & 0.001 & 0.015 & 0.001 & 15 & 5 \\
(5) Active fintech only & -1.110* & (0.478) & 0.028 & 0.142 & 0.052 & 27 & 2 \\
(6) PPML (exposure = clients) & -1.583*** & (0.242) & 0.000 &  &  & 29 & 4 \\
\bottomrule
\end{tabular}
}
\par\vspace{2pt}\parbox{\textwidth}{\footnotesize Dependent variable: ln(1 + fraud events per 100,000 retail clients) in specs (1)--(5); fraud event counts in spec (6) (Poisson pseudo-maximum likelihood, clients as exposure; coefficient is a semi-elasticity). All models include entity and month fixed effects; post = 2024m1 onward. Stars follow the permutation p-value (CR1 p-value in (6)). WCR: wild cluster bootstrap-t with Webb weights, null imposed (4,999 draws). Permutation: 20,000 random re-assignments of the fintech label across entities.}
\end{table}